\documentclass[11pt]{article}
% ======================================================================
%  examstyle.tex — EPFL-style exam layout for the weekly Time Series exams
%  Shared by every Exam_Week_NN(.tex / _Solutions.tex). Compiles with pdflatex.
% ======================================================================
\usepackage[utf8]{inputenc}
\usepackage[T1]{fontenc}
\usepackage[english]{babel}
\usepackage[a4paper,margin=2.4cm]{geometry}
\usepackage{amsmath,amssymb,amsthm,mathtools}
\usepackage{bm}
\usepackage{enumitem}
\usepackage{array}

\setlength{\parindent}{0pt}
\setlength{\parskip}{0.5em}
\emergencystretch=3em
\allowdisplaybreaks[1]

% ---------- math macros (same names as the rest of the project) ----------
\newcommand{\E}{\mathbb{E}}
\DeclareMathOperator{\Var}{Var}
\DeclareMathOperator{\Cov}{Cov}
\DeclareMathOperator{\Corr}{Corr}
\DeclareMathOperator*{\argmin}{arg\,min}
\DeclareMathOperator{\MSE}{MSE}
\DeclareMathOperator{\trace}{tr}
\newcommand{\Reals}{\mathbb{R}}
\newcommand{\Ints}{\mathbb{Z}}
\newcommand{\Comp}{\mathbb{C}}
\newcommand{\Nat}{\mathbb{N}}
\newcommand{\Prob}{\mathbb{P}}
\newcommand{\Tr}{^{\mathsf{T}}}
\newcommand{\conj}[1]{\overline{#1}}
\newcommand{\abs}[1]{\left\lvert #1 \right\rvert}
\newcommand{\norm}[1]{\left\lVert #1 \right\rVert}
\newcommand{\ind}{\mathbf{1}}
\newcommand{\eps}{\varepsilon}
\newcommand{\diff}{\nabla}
\newcommand{\acvs}{\gamma}
\newcommand{\acf}{\rho}
\newcommand{\pacf}{\alpha}
\newcommand{\sdf}{\mathcal{S}}
\newcommand{\Bsh}{B}
\newcommand{\iid}{\overset{\text{iid}}{\sim}}

\setlist[enumerate]{leftmargin=2em,topsep=2pt,itemsep=4pt}
\setlist[itemize]{leftmargin=1.6em,topsep=2pt,itemsep=2pt}

\newcounter{exo}

% ---------- exam cover header :  \examheader{exam title}{duration} ----------
\newcommand{\examheader}[2]{%
  \thispagestyle{empty}
  \begin{center}
    {\Large\bfseries Time Series}\\[3pt]
    Spring Semester 2026, EPFL\\
    \'Ecole Polytechnique F\'ed\'erale de Lausanne\\
    Prof.\ Dr.\ Sofia C.\ Olhede\\[7pt]
    \hrule height 0.8pt
    \vspace{9pt}
    {\large\bfseries Practice Exam --- #1}\\[4pt]
    {\normalsize Duration: #2 \quad$\cdot$\quad No calculator \quad$\cdot$\quad Answer on paper}\\
    \vspace{8pt}
    \hrule height 0.8pt
  \end{center}
  \vspace{14pt}
  \begin{tabular}{@{}p{8.2cm}@{}p{8cm}@{}}
    Name:\enspace\hrulefill & Surname:\enspace\hrulefill
  \end{tabular}\\[14pt]
  SCIPER (Student Card Number):\enspace\hrulefill
  \vspace{16pt}
}

% ---------- points table (fixed: 4 problems) ----------
\newcommand{\pointstable}{%
  \begin{center}
  \renewcommand{\arraystretch}{1.5}
  \begin{tabular}{|p{5cm}|p{4cm}|}
    \hline
    \textbf{Exercise} & \textbf{Points}\\\hline
    1 & \\\hline
    2 & \\\hline
    3 & \\\hline
    4 & \\\hline
    \textbf{Total} & \\\hline
  \end{tabular}
  \end{center}
}

% ---------- problem heading (exam: one problem per page) ----------
\newcommand{\problem}[1]{%
  \clearpage
  \stepcounter{exo}%
  \begin{center}\textbf{\large Problem \theexo\ (#1 points)}\end{center}
  \vspace{4pt}
}
% ---------- answer space: fill the statement page + one blank answer page ----------
% Put \answerpage at the END of each problem so every exercise gets its own page(s).
\newcommand{\answerpage}{\par\vfill\newpage\null\newpage}

% ---------- solutions document ----------
\newcommand{\solheader}[1]{%
  \begin{center}
    {\Large\bfseries Time Series --- Solutions}\\[3pt]
    {\large Practice Exam --- #1}\\[2pt]
    EPFL, MATH-342 \quad$\cdot$\quad Prof.\ S.\ C.\ Olhede\\[5pt]
    \hrule height 0.8pt
  \end{center}
  \vspace{10pt}
}
\newcommand{\solproblem}[1]{%
  \bigskip
  \stepcounter{exo}%
  {\large\bfseries Problem \theexo\ (#1 points)}\par\nobreak\vspace{2pt}
}
% Rappel de l'énoncé avant la solution (dans les corrigés)
\newcommand{\statementstart}{\par\vspace{1pt}{\bfseries\itshape Statement.}\par\nobreak\begingroup\small\itshape}
\newcommand{\statementend}{\par\endgroup\vspace{5pt}\hrule\vspace{6pt}{\bfseries Solution.}\par\nobreak\vspace{2pt}}

\begin{document}

\examheader{Week 6: Fourier Theory}{60 minutes}
\pointstable
\begin{center}\itshape Justify every answer. Throughout, $\eps_t$ is a mean-zero white noise of variance $\sigma^2$.\end{center}

% ---------------------------------------------------------------
\problem{5}
Throughout this problem $g\in L^1$ has continuous-time Fourier transform
\[
G(f)=\int_{-\infty}^{\infty} g(t)\,e^{-2\pi i f t}\,dt ,
\]
and for a derived function $h$ we write $H$ for its transform. We use the four elementary Fourier-transform properties (conjugation, scaling, shift, linearity).

\textbf{(a)} Prove the \emph{conjugation} and \emph{time-shift} rules: if $h(t)=\conj{g(t)}$ then $H(f)=\conj{G(-f)}$, and if $h(t)=g(t+\tau)$ then $H(f)=G(f)\,e^{2\pi i f\tau}$. \hfill(2 pts)

\textbf{(b)} Combining scaling and shifting, show that the \emph{modulated} signal $h(t)=g(\alpha t)\,e^{2\pi i f_0 t}$ (with $\alpha\neq0$, $f_0\in\Reals$ fixed) has transform
\[
H(f)=\frac{1}{\abs\alpha}\,G\!\left(\frac{f-f_0}{\alpha}\right).
\]
State in one sentence what each of the two operations (scaling by $\alpha$, modulation by $e^{2\pi i f_0 t}$) does to the spectrum. \hfill(2 pts)

\textbf{(c)} Use part (a) to show that if $g$ is \emph{real and even} (so $g(t)=\conj{g(t)}=g(-t)$), then its transform $G$ is real and even. \hfill(1 pt)
\answerpage

% ---------------------------------------------------------------
\problem{5}
Recall the continuous-time convolution of $g,h\in L^1$,
\[
(g*h)(t)=\int_{-\infty}^{\infty} g(t-y)\,h(y)\,dy ,
\]
and that upper-case letters denote Fourier transforms.

\textbf{(a)} Prove the \emph{convolution theorem}: if $u=g*h$ then $U=G\cdot H$. State explicitly the Fubini step and why the hypotheses $g,h\in L^1$ justify it. \hfill(3 pts)

\textbf{(b)} \emph{(Filtering.)} A signal $x\in L^1$ is passed through a linear filter with impulse response $a\in L^1$, producing $y=a*x$. Using part (a), give $Y$ in terms of $X$ and the transfer function $A$. If the filter is the running-average kernel $a(t)=\tfrac1{2T}\,\ind_{\{\abs t\le T\}}$, compute $A(f)$ in closed form and verify $A(0)=1$. \hfill(2 pts)
\answerpage

% ---------------------------------------------------------------
\problem{5}
This problem concerns the finite-data Fourier transform, the Dirichlet kernel, and aliasing. Take sampling interval $\Delta=1$ and observation times $t=0,1,\dots,n-1$.

\textbf{(a)} The transform of the finite rectangular window is the \textbf{Dirichlet kernel}
\[
D_n(f)=\sum_{t=0}^{n-1} e^{-2\pi i f t}.
\]
Sum this finite geometric series and show that
\[
D_n(f)=\frac{\sin(n\pi f)}{\sin(\pi f)}\;e^{-\pi i (n-1) f},
\qquad\text{so}\qquad
\abs{D_n(f)}=\left\lvert\frac{\sin(n\pi f)}{\sin(\pi f)}\right\rvert .
\]
\hfill(2 pts)

\textbf{(b)} Show that $\abs{D_n(f)}$ vanishes exactly at the nonzero \textbf{Fourier frequencies} $f_k=k/n$, $k=1,\dots,n-1$, while $D_n(0)=n$. (This is why the finite Fourier transform of a pure sinusoid sitting exactly on a Fourier frequency contributes nothing at the other Fourier frequencies.) \hfill(1.5 pts)

\textbf{(c)} \emph{(Aliasing / Nyquist.)} Let $g\in L^1$ have continuous-time Fourier transform $\mathcal G(f)$, and let $G$ denote the transform of the sampled sequence $\{g(t)\}_{t\in\Delta\Ints}$. Recall the aliasing identity
\[
G(f)=\sum_{k\in\Ints}\mathcal G\!\left(f+\frac{k}{\Delta}\right),
\qquad \abs f\le \frac{1}{2\Delta}.
\]
Suppose $\mathcal G$ is band-limited triangular,
$\mathcal G(f)=(1-\abs f)\,\ind_{\{\abs f\le 1\}}$.
For $\Delta=1$ (Nyquist frequency $\tfrac12$) compute the aliased transform $G(f)$ on $\abs f\le\tfrac12$, and state whether aliasing has occurred and why. \hfill(1.5 pts)
\answerpage

% ---------------------------------------------------------------
\problem{5}
\textbf{(Revision of Week 5: differencing \& SARIMA.)} Consider the signal-plus-noise model $X_t=\mu_t+Y_t$, where $\mu_t$ is a deterministic component and $\{Y_t\}$ is a zero-mean stationary process with autocovariance sequence $\acvs_\tau$. Recall $\diff=I-\Bsh$, $\diff_{(s)}=I-\Bsh^{s}$.

\textbf{(a)} Let $\mu_t=a+bt+s_t$ with a \emph{monthly} season $s_{t+12}=s_t$. Compute $\diff\,\diff_{(12)}X_t$ explicitly as a linear combination of the $Y$'s, verify its mean is $0$, and explain in one sentence why it is weakly stationary. \hfill(2 pts)

\textbf{(b)} Now let $\{X_t\}$ be the multiplicative seasonal process $\mathrm{SARMA}(0,1)\times(1,0)_{12}$,
\[
\bigl(1-\Phi\,\Bsh^{12}\bigr)X_t=\bigl(1-\theta\Bsh\bigr)\eps_t,
\qquad \abs\Phi<1 .
\]
By inverting the seasonal AR factor, obtain the $\mathrm{MA}(\infty)$ weights $\psi_j$ in $X_t=\sum_{j\ge0}\psi_j\eps_{t-j}$ (closed form for $\psi_{12l}$ and $\psi_{12l+1}$; all other $\psi_j=0$). \hfill(2 pts)

\textbf{(c)} Using $\acvs_\tau=\sigma^{2}\sum_{j\ge0}\psi_j\psi_{j+\tau}$, compute $\acvs_0$ and $\acvs_1$ in closed form, and evaluate $\acf_1=\acvs_1/\acvs_0$ for $\theta=\tfrac12$. \hfill(1 pt)
\answerpage

\end{document}
